\section{Escopo e limitações}
\label{sec:escopo-limitacoes}

O trabalho abrange o desenvolvimento do módulo de solicitações e as
funcionalidades necessárias ao gerenciamento das operações de alteração
cadastral e de classificação dos especialistas no sistema Banco de
Especialistas. Estão incluídos os mecanismos de registro, acompanhamento,
decisão e reversão das solicitações, bem como os recursos de persistência e
rastreabilidade associados a esse fluxo.

Para os fins deste trabalho, adotam-se as seguintes definições:

\begin{itemize}
  \item \textbf{Solicitação de complemento cadastral:} requisição para que o
  especialista seja submetido a nova avaliação, com a alteração de seu status
  para COMPLEMENTAÇÃO e o consequente recálculo de sua pontuação, conforme as
  regras de negócio aplicáveis.

  \item \textbf{Solicitação de mudança de classificação:} requisição de
  alteração da categoria do especialista entre as categorias existentes
  (estudante, colaborador externo e servidor), observadas as regras de negócio
  que impedem, por exemplo, a transição para a mesma categoria e a aprovação
  quando já existe retificação em curso.
\end{itemize}

Não fazem parte do escopo a redefinição dos critérios de pontuação dos
especialistas, a modificação do processo seletivo, a substituição do sistema de
avaliação existente e a implementação de notificações por e-mail. Também não é
contemplado o tratamento de múltiplas solicitações simultâneas para um mesmo
especialista.

O desenvolvimento foi conduzido por um único desenvolvedor, com restrições de
prazo e de experiência. As decisões de arquitetura e de implementação foram
fundamentadas em pesquisa, documentação técnica e experimentação, e validadas
por testes automatizados e entregas incrementais. Como consequência, algumas
escolhas privilegiam simplicidade e aderência ao sistema legado, em vez de uma
solução ótima para todos os cenários.